\section{Initial frontiers near a mark}
\label{sec:area_law_initial_star_frontiers}

Fix an arbitrary origin, finite endpoint set and valid residue choices
for the initial construction, with $C\ge2$ and
$k_0\ge50{,}000{,}000$. The estimate below is the boundary consequence
of the local mesh argument in \cite[Section~11]{OpenAI2026AreaLaw}.

% Source: 10-geometry.tex, geometry:initial-stars, lines 333--370,
% especially 352--359; prop:two-families, lines 154--177, 212--218
% and 299--323; revision adc7f1241b42e322a6451854ab7e4b4c146bf78a.

\begin{theorem}[A nearby initial frontier has an allowed direction]
    \label{thm:area_law_initial_frontier_near_mark}
    \lean{TNLean.PEPS.AreaLaw.Geometry.initialBirthRegion_frontier_near_mark_allowed}
    \leanok
    \uses{def:area_law_initial_birth_regions,def:area_law_belt_cell_marks,
        def:area_law_fine_layer_indices}
    Let $k\ge k_0$, let $z$ be an actual fine-cell index in layer $k$,
    and let $v$ be one of its nine marks. Put $t=2^{\ell_k}$.
    For any actual initial identifier $i$ and any
    $x\in\partial X_i^0$ satisfying $\|x-v\|_\infty<t/32$, the vector
    $x-v$ has one of the four allowed slopes. Thus $x$ lies on one
    of the four allowed lines through $v$.
\end{theorem}
\begin{proof}\leanok
    \uses{thm:area_law_primary_fine_cell_cover,def:area_law_fan_runs,
        thm:area_law_belt_marks_quarter_mesh,thm:area_law_square_frontier_side,
        thm:area_law_fan_frontier_primitives,
        thm:area_law_near_mark_quarter_mesh,
        thm:area_law_dummy_fine_scale_separation,
        thm:area_law_affine_mesh_line_clearance}
    The frontier of a finite union is contained in the union of the
    constituent frontiers. For a primary, its exact finite closed-cell
    cover therefore supplies a square frontier; for a fan run, its
    finite triangle support supplies an outer square side or an actual
    radial segment. The endpoints of each primitive segment are actual
    cell marks, and its slope is allowed by the triangle construction.

    Contact with a primitive segment supplies a point of its actual
    closed fine cell within $t/32$ of $v$. The nearby-cell estimate puts
    both segment endpoints and $v$ on the common quarter mesh. An
    allowed supporting line missing $v$ has distance at least $t/8$
    from $v$, a contradiction. Both points therefore lie on the same
    allowed line, giving the asserted direction.

    A dummy frontier point belongs to the closed dummy neighborhood.
    If $k>k_0$, the sixteen-fine-side separation contradicts its stated
    distance from $v$. Hence $k=k_0$. A constituent coarse-square
    frontier supplies a whole side, whose endpoints lie on the same
    quarter mesh by the established mesh inclusion and
    $\ell_k\le k$. The same line-clearance argument applies.
\end{proof}

This estimate concerns primitive boundary directions. Assigning actual
identifiers to the local sectors and identifying the active rays are
subsequent assertions.
